% Created 2026-09-15 mar 09:00
% Intended LaTeX compiler: pdflatex
\documentclass[11pt]{article}
\usepackage[utf8]{inputenc}
\usepackage[T1]{fontenc}
\usepackage{graphicx}
\usepackage{longtable}
\usepackage{wrapfig}
\usepackage{rotating}
\usepackage[normalem]{ulem}
\usepackage{amsmath}
\usepackage{amssymb}
\usepackage{capt-of}
\usepackage{hyperref}
\author{Marisol}
\date{2025}
\title{TP1 - Estructuras de Datos Avanzadas y Análisis Amortizado\\\medskip
\large Diseño de Algoritmos}
\hypersetup{
 pdfauthor={Marisol},
 pdftitle={TP1 - Estructuras de Datos Avanzadas y Análisis Amortizado},
 pdfkeywords={},
 pdfsubject={},
 pdfcreator={Emacs 29.3 (Org mode 9.6.15)}, 
 pdflang={Spanish}}
\begin{document}

\maketitle
\setcounter{tocdepth}{2}
\tableofcontents


\section{Datos del trabajo}
\label{sec:org3fc47ad}

\begin{itemize}
\item \textbf{Materia:} Diseño de Algoritmos
\item \textbf{Trabajo práctico:} TP1
\item \textbf{Tema:} Estructuras de Datos Avanzadas y Análisis Amortizado
\end{itemize}

\section{Ejercicio I-1: Montículo Binomial}
\label{sec:orgef7dd99}

\subsection{A) Diseño UML de clases}
\label{sec:orgc4b12c8}

\begin{quote}
Dibujar el diagrama UML de clases para implementar un montículo binomial dinámico.
\end{quote}

El diseño debe incluir las siguientes clases principales:

\begin{itemize}
\item \texttt{NodoBinomial}
\item \texttt{ArbolBinomial}
\item \texttt{HeapBinomial}

\item[{$\square$}] Definir las clases principales.
\item[{$\square$}] Identificar los atributos y métodos.
\end{itemize}



\subsubsection{Diagrama UML}
\label{sec:orgc26d299}

\begin{verbatim}
@startuml

class NodoBinomial {
  ' Completar atributos y métodos
}

class ArbolBinomial {
  ' Completar atributos y métodos
}

class HeapBinomial {
  ' Completar atributos y métodos
}

' Completar relaciones y multiplicidades

@enduml
\end{verbatim}

\subsubsection{Decisiones de diseño}
\label{sec:org079cd21}

\begin{itemize}
\item ¿Qué información almacena cada \texttt{NodoBinomial}?
\item ¿Cómo se representa la raíz de un árbol binomial?
\item ¿Cómo se enlazan los árboles binomiales dentro del heap?
\item ¿Qué atributo permite localizar rápidamente el mínimo?
\end{itemize}

\begin{verbatim}
Escribir aquí las decisiones de diseño.
\end{verbatim}

\subsection{B) Algoritmos}
\label{sec:orgb646ac8}

Desarrollar los algoritmos para las siguientes operaciones:

\begin{itemize}
\item \texttt{insertar(x)}
\item \texttt{unir(H1,H2)}
\item \texttt{buscarMin()}
\item \texttt{extraerMin()}
\item \texttt{disminuirClave(x,k)}
\item \texttt{eliminar(x)}
\end{itemize}

\subsubsection{\texttt{insertar(x)}}
\label{sec:org46ece6b}

\begin{verbatim}
Escribir aquí el pseudocódigo de insertar(x).
\end{verbatim}

\subsubsection{\texttt{unir(H1,H2)}}
\label{sec:org1221b18}

\begin{verbatim}
Escribir aquí el pseudocódigo de unir(H1,H2).
\end{verbatim}

\subsubsection{\texttt{buscarMin()}}
\label{sec:org1179a0b}

\begin{verbatim}
Escribir aquí el pseudocódigo de buscarMin().
\end{verbatim}

\subsubsection{\texttt{extraerMin()}}
\label{sec:org7ec6bd1}

\begin{verbatim}
Escribir aquí el pseudocódigo de extraerMin().
\end{verbatim}

\subsubsection{\texttt{disminuirClave(x,k)}}
\label{sec:org8964af9}

\begin{verbatim}
Escribir aquí el pseudocódigo de disminuirClave(x,k).
\end{verbatim}

\subsubsection{\texttt{eliminar(x)}}
\label{sec:org2ddf2f1}

\begin{verbatim}
Escribir aquí el pseudocódigo de eliminar(x).
\end{verbatim}

\subsection{C) Análisis de eficiencia temporal}
\label{sec:org856e5bb}

Analizar la eficiencia temporal de cada operación en notación \texttt{O} y justificarla.

\begin{center}
\begin{tabular}{lll}
Operación & Complejidad & Justificación\\[0pt]
\hline
\texttt{insertar(x)} &  & \\[0pt]
\texttt{unir(H1,H2)} &  & \\[0pt]
\texttt{buscarMin()} &  & \\[0pt]
\texttt{extraerMin()} &  & \\[0pt]
\texttt{disminuirClave(x,k)} &  & \\[0pt]
\texttt{eliminar(x)} &  & \\[0pt]
\end{tabular}
\end{center}

\begin{verbatim}
Escribir aquí la justificación detallada de las complejidades.
\end{verbatim}

\section{Ejercicio I-2: Conjuntos Disjuntos}
\label{sec:org28ae0bc}

\subsection{A) Implementación con arreglos}
\label{sec:orgf517c13}

Implementar en pseudocódigo o Java las operaciones sobre conjuntos disjuntos usando una representación en arreglos:

\begin{itemize}
\item \texttt{buscar(x)} (\texttt{find})
\item \texttt{fusionar(a,b)} (\texttt{union})
\end{itemize}

\subsubsection{Representación utilizada}
\label{sec:org5726f8e}

\begin{verbatim}
Explicar qué representa cada posición del arreglo y cómo se inicializan los conjuntos.
\end{verbatim}

\subsubsection{Implementación en pseudocódigo}
\label{sec:org006d89d}

\begin{verbatim}
Escribir aquí la implementación de buscar(x) y fusionar(a,b).
\end{verbatim}

\subsubsection{Implementación en Java}
\label{sec:orgc894545}

\begin{verbatim}
public class ConjuntosDisjuntos {

    // Completar atributos.

    public ConjuntosDisjuntos(int n) {
        // Completar inicialización.
    }

    public int buscar(int x) {
        // Completar implementación.
        return 0;
    }

    public void fusionar(int a, int b) {
        // Completar implementación.
    }
}
\end{verbatim}

\subsection{B) Mejoras de eficiencia}
\label{sec:orga0b4196}

Explicar cómo cambia la eficiencia al aplicar las siguientes técnicas:

\subsubsection{Union by rank}
\label{sec:org0f6bac9}

\begin{quote}
Explicar cómo se decide qué raíz se convierte en hija de la otra y por qué esto evita que los árboles crezcan demasiado.
\end{quote}

\begin{verbatim}
Escribir aquí la explicación de union by rank.
\end{verbatim}

\subsubsection{Path compression}
\label{sec:org72484a3}

\begin{quote}
Explicar cómo se modifican los enlaces durante la operación \texttt{buscar(x)} para acercar los nodos a la raíz.
\end{quote}

\begin{verbatim}
Escribir aquí la explicación de path compression.
\end{verbatim}

\subsection{C) Complejidad y justificación}
\label{sec:orgaa3bc6d}

Calcular la complejidad de las operaciones en cada caso y justificar.

\begin{center}
\begin{tabular}{llll}
Variante & Complejidad de \texttt{buscar} & Complejidad de \texttt{fusionar} & Justificación\\[0pt]
\hline
Representación básica &  &  & \\[0pt]
Union by rank &  &  & \\[0pt]
Path compression &  &  & \\[0pt]
Union by rank + path compression &  &  & \\[0pt]
\end{tabular}
\end{center}

\begin{verbatim}
Escribir aquí el análisis completo. Indicar, si corresponde, la complejidad amortizada.
\end{verbatim}

\section{Ejercicio I-3: Árbol TRIE - Diccionario de sinónimos}
\label{sec:org5b111f7}

\subsection{A) Diseño e implementación}
\label{sec:org9ebff28}

Diseñar e implementar en Java un algoritmo que permita:

\begin{itemize}
\item[{$\square$}] Almacenar un diccionario de sinónimos en un TRIE.
\item[{$\square$}] Agregar un sinónimo a una palabra existente.
\item[{$\square$}] Mostrar todos los sinónimos de una palabra dada.
\item[{$\square$}] Listar todas las palabras del diccionario sin mostrar sus sinónimos.
\end{itemize}

\subsubsection{Estructura de datos}
\label{sec:org2476f08}

Describir qué información contiene cada nodo del TRIE y cómo se representa la lista de sinónimos asociada a una palabra.

\begin{verbatim}
Escribir aquí la descripción de la estructura.
\end{verbatim}

\subsubsection{Implementación en Java}
\label{sec:orgd2c22b1}

\begin{verbatim}
import java.util.*;

public class TrieSinonimos {

    private static class Nodo {
        // Completar atributos del nodo.
    }

    private final Nodo raiz;

    public TrieSinonimos() {
        // Completar inicialización.
        raiz = null;
    }

    public void insertarPalabra(String palabra) {
        // Completar implementación.
    }

    public void agregarSinonimo(String palabra, String sinonimo) {
        // Completar implementación.
    }

    public List<String> mostrarSinonimos(String palabra) {
        // Completar implementación.
        return new ArrayList<>();
    }

    public List<String> listarPalabras() {
        // Completar implementación.
        return new ArrayList<>();
    }
}
\end{verbatim}

\subsubsection{Ejemplo de uso}
\label{sec:org255a8d6}

\begin{verbatim}
// Completar con un ejemplo de creación del diccionario,
// inserción de palabras, carga de sinónimos y consultas.
\end{verbatim}

\subsection{B) Análisis de eficiencia}
\label{sec:org8e856ed}

Analizar la eficiencia temporal de las operaciones \texttt{insertar}, \texttt{buscar} y \texttt{listar}.

Sea \texttt{L} la longitud de la palabra y sea \texttt{N} la cantidad de palabras almacenadas en el diccionario.

\begin{center}
\begin{tabular}{lll}
Operación & Complejidad & Justificación\\[0pt]
\hline
Insertar una palabra &  & \\[0pt]
Buscar una palabra &  & \\[0pt]
Agregar un sinónimo &  & \\[0pt]
Mostrar sinónimos &  & \\[0pt]
Listar todas las palabras &  & \\[0pt]
\end{tabular}
\end{center}

\begin{verbatim}
Escribir aquí la justificación de cada complejidad.
\end{verbatim}

\section{Ejercicio I-4: Análisis Amortizado}
\label{sec:org468fb29}

\subsection{A) Contador binario extendido}
\label{sec:orgd60f424}

Demostrar que, si se incluye la operación de decrementar una unidad en el contador binario visto en clase, una secuencia de \texttt{n} operaciones costará a lo sumo:

\[
\Theta(n \cdot k)
\]

donde \texttt{k} es la cantidad de dígitos del número binario.

\subsubsection{Modelo del contador}
\label{sec:org581b7f7}

Describir cómo se representa el contador y cómo funcionan las operaciones \texttt{incrementar} y \texttt{decrementar}.

\begin{verbatim}
Escribir aquí el modelo del contador.
\end{verbatim}

\subsubsection{Costo de incrementar}
\label{sec:org0f76416}

\begin{verbatim}
Analizar cuántos bits pueden cambiar durante un incremento.
\end{verbatim}

\subsubsection{Costo de decrementar}
\label{sec:org6d75c5d}

\begin{verbatim}
Analizar cuántos bits pueden cambiar durante un decremento.
\end{verbatim}

\subsubsection{Demostración para una secuencia de n operaciones}
\label{sec:orgdc2b95d}

\begin{verbatim}
Escribir aquí la demostración y explicar por qué el costo total está acotado por Θ(n·k).
\end{verbatim}

\subsection{B) Secuencia de operaciones con costos especiales}
\label{sec:org300614e}

\subsubsection{Enunciado}
\label{sec:orgf0b4909}

Se realiza una secuencia de \texttt{n} operaciones sobre una estructura de datos dada. La operación \texttt{i}-ésima cuesta:

\begin{itemize}
\item \texttt{costo(i) = i} si \texttt{i} es una potencia de 2.
\item \texttt{costo(i) = 1} en cualquier otro caso.
\end{itemize}

\subsubsection{Método agregado}
\label{sec:org1fcb5b1}

Usando el método agregado, calcular el costo amortizado de cada operación.

\begin{verbatim}
Escribir aquí el cálculo del costo total de la secuencia.
\end{verbatim}

\subsubsection{Costo total}
\label{sec:org6c0d6d2}

La suma de los costos puede expresarse como:

\[
T(n) = \sum_{i=1}^{n} costo(i)
\]

\begin{verbatim}
Desarrollar y acotar la suma de los costos de las operaciones que son potencias de 2.
\end{verbatim}

\subsubsection{Costo amortizado por operación}
\label{sec:org8764d14}

El costo amortizado promedio de una operación es:

\[
\widehat{c} = \frac{T(n)}{n}
\]

\begin{verbatim}
Completar el cálculo y expresar el resultado asintótico.
\end{verbatim}

\section{Conclusiones}
\label{sec:orgc0ef1f4}

\begin{verbatim}
Comparar las estructuras estudiadas y resumir qué técnica resulta más importante
en cada caso.
\end{verbatim}

\section{Bibliografía y fuentes}
\label{sec:org864f144}

\begin{itemize}
\item Material de la cátedra: TP1-1 - Análisis Amortizado - Diseño de Algoritmos, 2025.
\item Agregar aquí los libros, apuntes, artículos o páginas consultadas.
\end{itemize}

\section{Control de entrega}
\label{sec:org2b6f2ac}

\begin{itemize}
\item[{$\square$}] Ejercicio I-1 completo.
\item[{$\square$}] Diagrama UML incluido.
\item[{$\square$}] Algoritmos del montículo binomial incluidos.
\item[{$\square$}] Complejidades justificadas.
\item[{$\square$}] Ejercicio I-2 completo.
\item[{$\square$}] Union by rank explicado.
\item[{$\square$}] Path compression explicado.
\item[{$\square$}] Complejidades justificadas.
\item[{$\square$}] Ejercicio I-3 completo.
\item[{$\square$}] Código Java revisado y compilado.
\item[{$\square$}] Complejidades del TRIE justificadas.
\item[{$\square$}] Ejercicio I-4 completo.
\item[{$\square$}] Demostraciones revisadas.
\item[{$\square$}] Conclusiones escritas.
\item[{$\square$}] Bibliografía agregada.
\item[{$\square$}] Archivo exportado al formato solicitado por la cátedra.
\end{itemize}

\section{Notas}
\label{sec:org5bfea1f}

\begin{verbatim}
Escribir aquí dudas, observaciones o temas para consultar.
\end{verbatim}
\end{document}
